% Problem-section file following ../_template.tex.
\section{Entanglement of formation of generalized Bell-diagonal states}
\label{sec:problem-56}

\paragraph{Problem.}
For every local dimension $d\geq3$, determine the entanglement of formation
of an arbitrary Weyl--Bell-diagonal state.  Let
$\omega_d:=\exp(2\pi i/d)$ and define the generalized Pauli operators and
their associated Bell basis by
\begin{equation}
  X\lvert j\rangle:=\lvert j+1\!\!\pmod d\rangle,
  \qquad
  Z\lvert j\rangle:=\omega_d^j\lvert j\rangle,
  \qquad
  \lvert\Phi_{a,b}\rangle
  :=(I\otimes X^aZ^b)\lvert\Phi_d\rangle,
  \qquad
  \lvert\Phi_d\rangle:=\frac1{\sqrt d}\sum_{j=0}^{d-1}\lvert j,j\rangle,
  \label{eq:p56-weyl-bell-basis}
\end{equation}
where $a,b\in\mathbb Z_d$.  In this problem, ``Pauli-diagonal'' means
diagonal in the generalized Bell basis in
Eq.~\eqref{eq:p56-weyl-bell-basis}.  Thus the state is
\begin{equation}
  \rho_{\mathbf p}
  :=\sum_{a,b\in\mathbb Z_d}p_{a,b}
     \lvert\Phi_{a,b}\rangle\!\langle\Phi_{a,b}\rvert,
  \qquad
  p_{a,b}\geq0,
  \qquad
  \sum_{a,b\in\mathbb Z_d}p_{a,b}=1.
  \label{eq:p56-weyl-bell-diagonal-state}
\end{equation}
For every probability array $\mathbf p$ in
Eq.~\eqref{eq:p56-weyl-bell-diagonal-state}, determine an evaluable exact
formula and an optimal pure-state ensemble for
\begin{equation}
  E_F(\rho_{\mathbf p})
  :=\inf_{\rho_{\mathbf p}=\sum_i q_i
                    \lvert\psi_i\rangle\!\langle\psi_i\rvert}
       \sum_i q_i\,
       S\!\left(\operatorname{Tr}_B
          \lvert\psi_i\rangle\!\langle\psi_i\rvert\right),
  \qquad
  S(\sigma):=-\operatorname{Tr}(\sigma\log_2\sigma).
  \label{eq:p56-eof}
\end{equation}
The infimum in Eq.~\eqref{eq:p56-eof} is over finite pure-state ensembles
with $q_i\geq0$ and $\sum_iq_i=1$.

\subsection*{Status}

Unsolved

\subsection*{Source}

The problem is implicit in Vollbrecht and Werner's finite-Weyl symmetry
construction and general convex-roof reduction
\sourcecite{ref:p56-vollbrecht-werner}{VW01}.  Terhal and Vollbrecht's
solution of the isotropic subfamily isolates a one-parameter slice rather
than the full generalized Bell simplex
\sourcecite{ref:p56-terhal-vollbrecht}{TV00}.

\subsection*{Progress}

\begin{itemize}
  \item For $d=2$, Wootters determined $E_F$ for every two-qubit state and
  therefore for every ordinary Bell-diagonal state
  \sourcecite{ref:p56-wootters}{Woo98}.  This formula relies on the
  two-qubit concurrence and does not extend to the general $d\geq3$
  probability array in Eq.~\eqref{eq:p56-weyl-bell-diagonal-state}.

  \item The finite-Weyl twirl and the symmetry method of Vollbrecht and
  Werner reduce the full convex roof to a phase optimization followed by a
  convexification.  More precisely, define
  \begin{equation}
    \varepsilon_d(\mathbf p)
    :=\min_{\boldsymbol\theta\in[0,2\pi)^{d^2}}
       S\!\left(\operatorname{Tr}_B
          \lvert\psi_{\mathbf p,\boldsymbol\theta}\rangle
          \!\langle\psi_{\mathbf p,\boldsymbol\theta}\rvert\right),
    \qquad
    \lvert\psi_{\mathbf p,\boldsymbol\theta}\rangle
    :=\sum_{a,b\in\mathbb Z_d}
       \sqrt{p_{a,b}}e^{i\theta_{a,b}}\lvert\Phi_{a,b}\rangle.
    \label{eq:p56-symmetry-reduction}
  \end{equation}
  Their theorem gives
  \begin{equation}
    E_F(\rho_{\mathbf p})
    =\operatorname{co}\varepsilon_d(\mathbf p),
    \label{eq:p56-convexification}
  \end{equation}
  where $\operatorname{co}$ denotes the lower convex envelope on the
  probability simplex.  Equations~\eqref{eq:p56-symmetry-reduction} and
  \eqref{eq:p56-convexification} are an exact reduction, but the phase
  minimum and its convex envelope remain unknown for a general
  $\mathbf p$ \sourcecite{ref:p56-vollbrecht-werner}{VW01}.

  \item For the isotropic probability array
  \begin{equation}
    p_{0,0}=F,
    \qquad
    p_{a,b}=\frac{1-F}{d^2-1}\quad((a,b)\neq(0,0)),
    \qquad 0\leq F\leq1,
    \label{eq:p56-isotropic-array}
  \end{equation}
  Terhal and Vollbrecht reduced the answer to the convex hull of an explicit
  function \sourcecite{ref:p56-terhal-vollbrecht}{TV00}; Fei and Li-Jost
  subsequently proved the conjectured convex-hull shape in every dimension
  \sourcecite{ref:p56-fei-li-jost}{FLJ06}.  Hence the states in
  Eq.~\eqref{eq:p56-isotropic-array} are solved exactly, but they form only a
  one-dimensional subfamily of the $(d^2-1)$-dimensional simplex.

  \item A 2023 analysis of the same Weyl--Bell simplex classified Bell-diagonal
  qutrit and ququart states using analytical and numerical separability
  criteria, but left $22.6\%$ of its sampled PPT ququart states unclassified
  as separable or bound entangled
  \sourcecite{ref:p56-popp-hiesmayr}{PH23}.  Because $E_F(\rho)=0$ exactly for
  separable states, this unresolved zero-versus-positive boundary is already
  an obstruction to a general exact formula; that work does not evaluate
  $E_F$ on the unclassified states.
\end{itemize}

\subsection*{References}

\begin{enumerate}[label={},leftmargin=4.8em,labelsep=0.5em]
  \footnotesize
  \item[\textup{[VW01]}]\label{ref:p56-vollbrecht-werner}
  K. G. H. Vollbrecht and R. F. Werner,
  ``Entanglement Measures under Symmetry,''
  \emph{Physical Review A} \textbf{64}, 062307 (2001).
  \href{https://doi.org/10.1103/PhysRevA.64.062307}%
       {doi:10.1103/PhysRevA.64.062307};
  \href{https://arxiv.org/abs/quant-ph/0010095}%
       {arXiv:quant-ph/0010095}.

  \item[\textup{[TV00]}]\label{ref:p56-terhal-vollbrecht}
  B. M. Terhal and K. G. H. Vollbrecht,
  ``Entanglement of Formation for Isotropic States,''
  \emph{Physical Review Letters} \textbf{85}, 2625--2628 (2000).
  \href{https://doi.org/10.1103/PhysRevLett.85.2625}%
       {doi:10.1103/PhysRevLett.85.2625};
  \href{https://arxiv.org/abs/quant-ph/0005062}%
       {arXiv:quant-ph/0005062}.

  \item[\textup{[Woo98]}]\label{ref:p56-wootters}
  W. K. Wootters,
  ``Entanglement of Formation of an Arbitrary State of Two Qubits,''
  \emph{Physical Review Letters} \textbf{80}, 2245--2248 (1998).
  \href{https://doi.org/10.1103/PhysRevLett.80.2245}%
       {doi:10.1103/PhysRevLett.80.2245};
  \href{https://arxiv.org/abs/quant-ph/9709029}%
       {arXiv:quant-ph/9709029}.

  \item[\textup{[FLJ06]}]\label{ref:p56-fei-li-jost}
  S.-M. Fei and X. Li-Jost,
  ``$R$ Function Related to Entanglement of Formation,''
  \emph{Physical Review A} \textbf{73}, 024302 (2006).
  \href{https://doi.org/10.1103/PhysRevA.73.024302}%
       {doi:10.1103/PhysRevA.73.024302};
  \href{https://arxiv.org/abs/quant-ph/0602137}%
       {arXiv:quant-ph/0602137}.

  \item[\textup{[PH23]}]\label{ref:p56-popp-hiesmayr}
  C. Popp and B. C. Hiesmayr,
  ``Comparing Bound Entanglement of Bell Diagonal Pairs of Qutrits and
  Ququarts,'' \emph{Scientific Reports} \textbf{13}, 2037 (2023).
  \href{https://doi.org/10.1038/s41598-023-29211-w}%
       {doi:10.1038/s41598-023-29211-w};
  \href{https://arxiv.org/abs/2209.15267}{arXiv:2209.15267}.
\end{enumerate}

\subsection*{Comment}

The unresolved task is to evaluate the phase minimum in
Eq.~\eqref{eq:p56-symmetry-reduction}, determine its lower convex envelope in
Eq.~\eqref{eq:p56-convexification}, and construct optimal ensembles for
arbitrary $\mathbf p$ when $d\geq3$.  Problem~7 instead asks for the
regularized entanglement cost of qubit Bell-diagonal states.

\subsection*{Tag}

Bell-diagonal states; Entanglement measures; Qudit systems; Convex optimization

\subsection*{ID}

\texttt{op\_a3a8680c50800797}
